\subsection{交换环上两个模的张量积}

设 C 为一个交换环；对于每个 C-模 E，E 上的模结构与自身相容（§ 1，第 14 号）。若 E 和 F 是两个 C-模，则第 4 号的讨论使我们能够在张量积 E ⊗c F 上定义两个 C-模结构，分别使得 γ(x ⊗ y) = (γx) ⊗ y 以及 γ(x ⊗ y) = x ⊗ (γy)；但由于根据第 1 号定义 1，在此情形下 (γx) ⊗ y = x ⊗ (γy)，这两个结构是相同的。此后当我们把 E ⊗c F 作为 C-模来讨论时，除非另有说明，均指按此方式定义的结构。典范同构

\[
\sigma : \mathbf {F} \otimes_ {\mathbf {c}} \mathbf {E} \rightarrow \mathbf {E} \otimes_ {\mathbf {c}} \mathbf {F}
\]

（第 1 号，命题 1 的推论 2）则是一个 C-模同构。

由这一定义可知，若 \((a_{\lambda})_{\lambda\in\mathbb{L}}\) （相应地 \((b_{\mu})_{\mu\in\mathbb{M}}\) ）是 C-模 E（相应地 F）的一个生成系，则 \((a_{\lambda}\otimes b_{\mu})\) 是 C-模 \(E\otimes_{C}F\) 的一个生成系；特别地，若 E 和 F 是有限生成的 C-模，则 \(E\otimes_{C}F\) 亦然。

对于每个 C-模 G，那些 Z-双线性映射 \(f\) （从 \(\mathbf{E} \times \mathbf{F}\) 到 \(\mathbf{G}\) ），若它们满足

\[
f (\gamma x, y) = f (x, \gamma y) = \gamma f (x, y) \quad \text { for } x \in \mathbf {E}, y \in \mathbf {F}, \gamma \in \mathbf {C}\tag{9}
\]

于是称为 C-双线性映射，并构成一个 C-模，记作 \(\mathcal{L}_2(\mathrm{E},\mathrm{F};\mathrm{G})\) ；命题 3（第4号）定义了一个典范的 C-模同构（参见 §1，第 14 号，注记 1）。

\[
\mathcal {L} _ {2} (\mathrm{E}, \mathrm{F}; \mathrm{G}) \rightarrow \operatorname{Hom} _ {\mathbf {c}} (\mathrm{E} \otimes_ {\mathbf {c}} \mathrm{F}, \mathrm{G}).\tag{10}
\]

设 \(E'\) , \(F'\) 是两个 C-模，且 \(u: E \to E'\) , \(v: F \to F'\) 是两个 C-线性映射；则（第4号） \(u \otimes v\) 是 \(E \otimes_{\mathbb{C}} F\) 到 \(E' \otimes_{\mathbb{C}} F'\) 的 C-线性映射。进一步，立即可以看出 \((u, v) \mapsto u \otimes v\) 是 \(\operatorname{Hom}_{\mathbb{C}}(\mathbf{E}, \mathbf{E}') \times \operatorname{Hom}_{\mathbb{C}}(\mathbf{F}, \mathbf{F}')\) 到 \(\operatorname{Hom}_{\mathbb{C}}(\mathbf{E} \otimes_{\mathbb{C}} \mathbf{F}, \mathbf{E}' \otimes_{\mathbb{C}} \mathbf{F}')\) 的 C-双线性映射；因此典范地对应于它的一个 C-线性映射，称为典范映射：

\[
\operatorname{Hom} _ {\mathbf {c}} (\mathbf {E}, \mathbf {E} ^ {\prime}) \otimes_ {\mathbf {c}} \operatorname{Hom} _ {\mathbf {c}} (\mathbf {F}, \mathbf {F} ^ {\prime}) \to \operatorname{Hom} _ {\mathbf {c}} (\mathbf {E} \otimes_ {\mathbf {c}} \mathbf {F}, \mathbf {E} ^ {\prime} \otimes_ {\mathbf {c}} \mathbf {F} ^ {\prime})\tag{11}
\]

它使每个元素 \(u \otimes v\) （它属于张量积

\[
\operatorname{Hom} _ {\mathbf {c}} (\mathbf {E}, \mathbf {E} ^ {\prime}) \otimes_ {\mathbf {c}} \operatorname{Hom} _ {\mathbf {c}} (\mathbf {F}, \mathbf {F} ^ {\prime})
\]

）对应到线性映射 \(u \otimes v\) 。注意，典范映射（11）不一定是单射也不是满射（§4，习题2）。

注记. (1) 设 A, B 为两个交换环， \(\rho:B\to A\) 为一个环同态，且 E 与 F 为两个 A-模；则第 3 号中的典范映射 (6) 是一个 B-线性映射

\[
\rho_ {*} (\mathbf {E}) \otimes \rho_ {*} (\mathbf {F}) \rightarrow \rho_ {*} (\mathbf {E} \otimes_ {\mathbf {A}} \mathbf {F}).\tag{12}
\]

\begin{enumerate}
\def\labelenumi{(\arabic{enumi})}
\setcounter{enumi}{1}
\tightlist
\item
  本号中所述内容可推广到如下情形：设 E 为右 A-模，F 为左 A-模，C 为交换环，且 \(\rho:C\to A\) 为 C 到 A 的同态，使得 \(\rho(\mathbf{C})\) 包含在 A 的中心内（参见 III，§1，第 3 号）。于是我们可以考虑 C-模 \(\rho_{*}(\mathbf{E})\) 和 \(\rho_{*}(\mathbf{F})\) ，并且关于 \(\rho\) 的假设蕴含这些 C-模结构分别与 E 和 F 上的 A-模结构相容（§1，第 14 号）。张量积 \(E\otimes_{A}F\) 因此（根据第4号）被赋予两个 C-模结构，使得 \(\gamma(x\otimes y)=(x\rho(\gamma))\otimes y\) 和
\end{enumerate}

\[
\gamma (x \otimes y) = x \otimes (\rho (\gamma) y)
\]

分别对 \(\gamma\in\mathbf{C}\) , \(x\in\mathbf{E}\) , \(y\in\mathbf{F}\) 成立，并且定义 1（第 1 号）还表明这两个结构是相同的。若 \(\mathbf{E}'\) （相应地 \(\mathbf{F}'\) ）是一个右（相应地，左）A-模，且 \(u:\mathbf{E}\to\mathbf{E}'\) , \(v:\mathbf{F}\to\mathbf{F}'\) 是两个 A-线性映射，则 \(u\otimes v:\mathbf{E}\otimes_{\mathbf{A}}\mathbf{F}\to\mathbf{E}'\otimes_{\mathbf{A}}\mathbf{F}'\) 对于刚刚定义的 C-模结构而言是 C-线性的；映射 \((u,v)\mapsto u\otimes v\) :

\[
\operatorname{Hom} _ {\mathbf {A}} (\mathrm{E}, \mathrm{E} ^ {\prime}) \times \operatorname{Hom} _ {\mathbf {A}} (\mathrm{F}, \mathrm{F} ^ {\prime}) \rightarrow \operatorname{Hom} _ {\mathbf {C}} (\mathrm{E} \otimes_ {\mathbf {A}} \mathrm{F}, \mathrm{E} ^ {\prime} \otimes_ {\mathbf {A}} \mathrm{F} ^ {\prime})
\]

是 C-双线性的（对于 \(\mathrm{Hom}_{\mathbf{A}}(\mathbf{E},\mathbf{E}^{\prime})\) 和 \(\mathrm{Hom}_{\mathbf{A}}(\mathbf{F},\mathbf{F}^{\prime})\) 上的

在 §1，第 14 号，注记 1 中定义的 C-模结构），由此我们还导出一个 C-线性映射，称为典范映射

\[
\operatorname{Hom} _ {\mathbf {A}} (\mathbf {E}, \mathbf {E} ^ {\prime}) \otimes_ {\mathbf {C}} \operatorname{Hom} _ {\mathbf {A}} (\mathbf {F}, \mathbf {F} ^ {\prime}) \rightarrow \operatorname{Hom} _ {\mathbf {C}} (\mathbf {E} \otimes_ {\mathbf {A}} \mathbf {F}, \mathbf {E} ^ {\prime} \otimes_ {\mathbf {A}} \mathbf {F} ^ {\prime}).\tag{13}
\]

\begin{enumerate}
\def\labelenumi{(\arabic{enumi})}
\setcounter{enumi}{2}
\tightlist
\item
  设 \(\mathbf{A}\) 是一个整环， \(\mathbf{K}\) 为其分式域。若 \(\mathbf{E}\) 和 \(\mathbf{F}\) 是两个向量 \(\mathbf{K}\) -空间，则典范映射
\end{enumerate}

\[
\left(\mathbf {E} _ {[ \mathbf {A} ]}\right) \otimes_ {\mathbf {A}} \left(\mathbf {F} _ {[ \mathbf {A} ]}\right)\rightarrow \mathbf {E} \otimes_ {\mathbf {K}} \mathbf {F}
\]

（第 3 号及 §1，第 13 号）是双射的。只需（第4号）证明：若 \(f\) 是 \(\mathbf{E} \times \mathbf{F}\) 到向量 K-空间 \(\mathbf{G}\) 的一个 A-双线性映射，则 \(f\) 也是 K-双线性的。现在，对所有 \(\alpha \neq 0\) （在 \(\mathbf{A}\) 中）：

\[
\alpha f (\alpha^ {- 1} x, y) = f (x, y) = \alpha f (x, \alpha^ {- 1} y)
\]

，因此

\[
f (\alpha^ {- 1} x, y) = f (x, \alpha^ {- 1} y) = \alpha^ {- 1} f (x, y)
\]

，因为 G 是一个向量 K-空间。

